\documentclass[tikz,border=6pt]{standalone}
\usepackage{peabra}
\begin{document}
\begin{tikzpicture}[x=1cm,y=1cm,
  gy/.style={sotetdoboz=30mm, font=\footnotesize\bfseries},
  t/.style={doboz=30mm, font=\footnotesize\bfseries, minimum height=9mm},
  p/.style={font=\scriptsize, text=peSotet, align=left, text width=30mm, anchor=north}]
  \node[gy] (m) at (0,2.2) {modellek};
  \foreach \x/\n/\c/\p in {
     -7.2/f/{fizikai modellek}/{makett, skálamodell, prototípus, próbapad; a valódi fizikai jelenséget hordozza},
     -3.6/l/{fogalmi, leíró modellek}/{kontextusdiagram, funkcióhierarchia, folyamatábra, SysML-diagramok; szerkezetet és viselkedést ír le},
     0/a/{matematikai, analitikus modellek}/{egyenletek, képletek, optimalizálás; statikus vagy dinamikus, determinisztikus vagy sztochasztikus},
     3.6/s/{szimulációs modellek}/{diszkrét eseményes, rendszerdinamikai, ágensalapú, Monte-Carlo; a viselkedést időben játssza le},
     7.2/d/{adatvezérelt modellek}/{statisztikai és gépi tanulási modellek; a mért adatokból tanulja meg az összefüggést}} {
    \node[t] (\n) at (\x,0.4) {\c};
    \draw[nyilv] (m.south) -- (\n.north);
    \node[p] at (\x,-0.3) {\p};
  }
  % fehér-fekete doboz tengely
  \shade[left color=white, right color=peSotet] (-8.8,-3.4) rectangle (8.8,-3.0);
  \draw[peSzurke] (-8.8,-3.4) rectangle (8.8,-3.0);
  \node[font=\scriptsize, anchor=north west] at (-8.8,-3.5) {fehér doboz: a belső mechanizmus ismert (első elvek, fizikai törvények)};
  \node[font=\scriptsize, anchor=north east] at (8.8,-3.5) {fekete doboz: csak a bemenet–kimenet kapcsolat ismert (mért adatok)};
  \node[font=\scriptsize, anchor=north] at (0,-3.95) {szürke doboz: ismert szerkezet, adatokból becsült paraméterek};
\end{tikzpicture}
\end{document}
